你跑完一个模型，屏幕上出现一个数：斜率 2.42，或者三组均值 15.3、16.1、18.7。把这个数写进报告，只需要抄。把它写成一句可以被别人依赖的话——``每多开一小时空调，用电量增加 2.4 度，上下浮动不超过 0.3''，或者``这三种程序的耗电确有差别''——中间隔着一段路。这一章走的就是这段路。

路上要付的代价是假设。点估计几乎不要假设：给定设计矩阵 \(X\)，最小二乘把观测 \(y\) 投影到列空间上，拆成可解释的部分与正交残差，这是纯粹的几何，任何一批数字都能做。可一旦你要说``上下浮动不超过 0.3''，就得回答另一个问题：换一批数据，这个数会漂到哪里去。几何对此保持沉默，答案只能来自一个关于误差如何生成的概率模型。误差独立、方差齐性、小样本下正态——这三条一加，几何分解立刻升级成系数的抽样分布、\(t\) 检验和 \(F\) 检验。

所以这一章始终把两笔账分开记。第一笔账是估计，第二笔账是推断。分开记的好处在诊断时才显现：异方差只弄坏第二笔账，系数依然无偏；分组或时序造成的相关误差同样只弄坏第二笔账，但它把标准误压得过窄，是三种失效里最容易让人虚报显著的一种；共线性把某些系数的抽样分布摊得很宽，却不妨碍拟合值和预测。三种毛病若都笼统地叫``模型有问题''，就没法对症。

同一段路走第二遍，场景换成比较多组均值。总平方和拆成组间与组内，比值构成 \(F\) 统计量，这套记法看起来自成一门技术。拆开看，它是把组标签编成指示变量后的回归，平方和分解就是上一节那次正交投影换了个写法。把这层对应关系认下来，协变量、随机效应、不平衡设计里那几种平方和的分歧，都能用同一套语言处理。

\subsection{怎么读这一章}

先把线性回归读成一个带随机误差的条件均值模型，而不只是一条最优拟合线；每读到一个推断结论，回头问它用掉了哪几条假设。再看 ANOVA 如何把总变异拆成组间与组内，并留意它与回归的对应处。两篇都要同步检查独立性、方差结构和数据是怎么采来的，因为显著性的合法性完全寄生在这些地方。

\subsection{本章篇目}

以下篇目按概念依赖排列；若从具体问题进入，也可以先打开最接近当前困惑的一篇，再沿篇内链接回补。

\begin{itemize}
\tightlist
\item \hyperref[sec:11-Mathematical-Statistics/08-Regression-and-ANOVA/01]{``经典线性回归把系数估计接到抽样分布''}；
\item \hyperref[sec:11-Mathematical-Statistics/08-Regression-and-ANOVA/02]{``ANOVA 用平方和分解比较多组均值''}。
\end{itemize}

读完这两篇，希望你对``显著''这个词的态度会变得具体一些。小 \(p\) 值不说明效应大，不说明模型对，也不说明是自变量或分组造成了结果的变化；它只说明，在你写明的零假设和误差模型下，眼前这个统计量不太容易出现。这句话里的每一个限定词，都是前面那笔假设账里的一项。
